\section{Calidad y preparación de datos}

La preparación se ejecuta desde la fuente inmutable mediante
\path{scripts/build_processed_dataset.py}. La implementación en
\path{src/spotify_popularity/data/preparation.py} valida esquema, detecta filas inválidas, comprueba
atributos invariantes y consolida IDs.

\begin{table}[H]
\centering
\caption{Problemas, decisiones y riesgos residuales de preparación.}
\label{tab:decisiones-preparacion}
\small
\begin{tabularx}{\textwidth}{p{2.5cm}p{3.1cm}p{2.7cm}X p{2.7cm}}
\toprule
Problema & Evidencia & Decisión & Justificación & Riesgo residual \\
\midrule
Índice exportado & Secuencia \code{Unnamed: 0} & Retirar del derivado & No describe una canción y podría introducir orden artificial & Ninguno relevante \\
Fila inválida & Tres textos nulos y duración cero & Excluir una fila & No hay evidencia para imputar identidad ni duración & Posible pérdida mínima de cobertura \\
IDs repetidos & 16.641 IDs; 24.259 apariciones extra & Consolidar por \code{track_id} & Evita doble conteo y fuga entre particiones & La popularidad histórica se resume \\
Géneros múltiples & Una pista puede aparecer en varias etiquetas & Unir etiquetas únicas con \code{|} & Conserva información sin duplicar canciones & Requiere encoding multietiqueta \\
Popularidad variable & 720 IDs repetidos difieren en el objetivo & Usar mediana & Es determinista y robusta sin privilegiar el orden del archivo & No representa actualidad ni trayectoria \\
\bottomrule
\end{tabularx}
\end{table}

Elegir la primera fila habría vinculado el resultado al orden del CSV; promediar sería más sensible
a discrepancias extremas; eliminar todos los IDs repetidos descartaría información válida. La
mediana y la preservación de etiquetas equilibran trazabilidad, robustez y retención de datos. El
código se detiene si atributos declarados invariantes entran en conflicto, evitando consolidaciones
silenciosas.

El contrato posterior confirma 89.740 filas y 20 columnas, 89.740 IDs únicos, cero celdas nulas,
cero duraciones no positivas y géneros no vacíos. Estas comprobaciones validan la entrada al EDA,
pero no eliminan sesgos de origen ni garantizan representatividad.

\subsection{Comparación antes y después}

\begin{table}[H]
\centering
\caption{Impacto cuantitativo de consolidar \code{track_id}. Tabla generada por
\code{scripts/build_latex_tables.py}.}
\label{tab:antes-despues}
\resizebox{\textwidth}{!}{\begin{tabular}{lrrrrrr}
\toprule
 & Filas & Canciones unicas & Media & Mediana & Porcentaje igual a 0 & Porcentaje mayor o igual a 70 \\
dataset &  &  &  &  &  &  \\
\midrule
Original & 114000 & 89741 & 33.239 & 35.000 & 14.053 & 4.800 \\
Procesado & 89740 & 89740 & 33.198 & 33.000 & 10.468 & 3.480 \\
\bottomrule
\end{tabular}
}
\end{table}

La tabla~\ref{tab:antes-despues} muestra que las canciones únicas pasan de 89.741 a 89.740 por la
exclusión de la fila inválida, mientras las filas caen 21,28\%. La media apenas cambia (33,24 a
33,20), pero la mediana baja de 35 a 33, la masa en cero de 14,05\% a 10,47\% y el segmento de 70 o
más de 4,80\% a 3,48\%. Esto demuestra que la duplicación por género modificaba la ponderación de
la distribución: incluso si la media permanece cercana, otras características relevantes cambian.
